% !TEX root = task-1-10-1.tex
% Editable task bank. Append new BankTask blocks at the end of the document.
\documentclass[12pt,a4paper]{article}
\usepackage{physics-textbook}
\usepackage{graphicx,needspace}
\usetikzlibrary{decorations.pathmorphing}
\newif\ifShowSolutions
\ifdefined\TutorEdition
  \ShowSolutionstrue
\else
  \ShowSolutionsfalse % Change to \ShowSolutionstrue to display solutions.
\fi
\newcounter{banktask}
\newsavebox{\BankTaskBox}
\newcommand{\BankTaskContent}[3]{%
  \noindent\textbf{\thebanktask.}\quad #2\par
  \ifShowSolutions
    \Needspace{50mm}% Keep the solution heading with its opening text or diagram.
    #3\par
    \ifdefined\TutorEdition
      \ifcsname TutorSource#1\endcsname
        \nopagebreak\smallskip{\small\csname TutorSource#1\endcsname\par}%
      \fi
    \fi
  \fi
}
\newcommand{\BankTask}[3]{%
  \refstepcounter{banktask}\par\medskip
  \sbox{\BankTaskBox}{\begin{minipage}{\linewidth}%
    \BankTaskContent{#1}{#2}{#3}\end{minipage}}%
  \ifdim\dimexpr\ht\BankTaskBox+\dp\BankTaskBox\relax>\dimexpr\textheight-20mm\relax
    \Needspace{35mm}\BankTaskContent{#1}{#2}{#3}%
  \else
    \Needspace{\dimexpr\ht\BankTaskBox+\dp\BankTaskBox\relax}%
    \noindent\usebox{\BankTaskBox}\par
  \fi
  \medskip
}
\begin{document}
\BookSubsection{1.10.1}{Закон всемирного тяготения}


% Source sdamgia-301: https://phys-ege.sdamgia.ru/problem?id=301
\BankTask{sdamgia-301}{
Две планеты с одинаковыми массами обращаются по круговым орбитам вокруг звезды. У первой из них радиус орбиты вдвое больше, чем у второй. Каково отношение сил притяжения первой и второй планет к звезде $\frac{F_1}{F_2}$?
}{
\smallskip\noindent\textbf{Решение}\par
По закону Всемирного тяготения сила притяжения планеты к звезде обратно пропорциональна квадрату радиуса орбиты. Таким образом, в силу равенства масс отношение сил притяжения к звезде первой и второй планет обратно пропорционально отношению квадратов радиусов орбит: 

\[\frac{F_1}{F_2} = \frac{R_2 ^2}{R_1 ^2} .\]

 По условию, у первой планеты радиус орбиты вдвое больше, чем у второй, то есть $R_1=2R_2,$ а значит, 

\[\frac{F_1}{F_2} = \frac{R_2 ^2}{( 2R_2 ) ^2} = \frac{R_2 ^2}{4R_2 ^2} = \frac{1}{4} =0{,}25.\]

Ответ: $0{,}25$.
}

% Source sdamgia-332: https://phys-ege.sdamgia.ru/problem?id=332
\BankTask{sdamgia-332}{
Космонавт на Земле притягивается к ней с силой $700\,\text{Н}$. С какой приблизительно силой он будет притягиваться к Марсу, находясь на его поверхности, если радиус Марса в $2$ раза меньше, а масса — в $10$ раз меньше, чем у Земли? (В ньютонах.)
}{
\smallskip\noindent\textbf{Решение}\par
Согласно закону всемирного тяготения сила притяжения двух тел пропорциональна произведению их масс и обратно пропорциональна квадрату расстояния между ними. Таким образом, силы притяжения космонавта к Земле и Марсу равны $F_{\text{Земли}}= \frac{GM_{\text{Земли}}m}{R_{\text{Земли}} ^2}$ и $F_{\text{Марса}}= \frac{GM_{\text{Марса}}m}{R_{\text{Марса}} ^2}$ соответственно. Отсюда находим приблизительно силу притяжения космонавта к Марсу: 

\[F_{\text{Марса}}=F_{\text{Земли}} \cdot  \frac{M_{\text{Марса}}}{M_{\text{Земли}}} \cdot  \frac{R_{\text{Земли}} ^2}{R_{\text{Марса}} ^2} =700\,\text{Н} \cdot  \frac{1}{10} \cdot  2 ^2 =280\,\text{Н}.\]

Ответ: $280$.
}

% Source sdamgia-6722: https://phys-ege.sdamgia.ru/problem?id=6722
\BankTask{sdamgia-6722}{
Расстояние от спутника до поверхности Земли равно радиусу Земли. Во сколько раз уменьшится сила притяжения спутника к Земле, если расстояние от него до поверхности Земли станет равным трем радиусам Земли?
}{
\smallskip\noindent\textbf{Решение}\par
Согласно закону всемирного тяготения, сила, с которой притягиваются точечное тело и сферически симметричное тело: $F=G \frac{m_1m_2}{r ^2} ,$ где r — это расстояние между центрами тел. В первом случае $F_1=G \frac{m_1m_2}{( R + R ) ^2} =G \frac{m_1m_2}{4R ^2} .$ Во втором случае $F_2=G \frac{m_1m_2}{( R + 3R ) ^2} =G \frac{m_1m_2}{16R ^2} ,$ то есть сила уменьшится в $4$ раза.

 

Ответ: $4$.
}

% Source sdamgia-6884: https://phys-ege.sdamgia.ru/problem?id=6884
\BankTask{sdamgia-6884}{
Сила притяжения Земли к Солнцу в $22{,}5$ раза больше, чем сила притяжения Марса к Солнцу. Во сколько раз расстояние между Марсом и Солнцем больше расстояния между Землей и Солнцем, если масса Земли в $10$ раз больше массы Марса?
}{
\smallskip\noindent\textbf{Решение}\par
Согласно закону всемирного тяготения сила притяжения между двумя материальными точками прямопропорциональна произведению их масс и обратнопропорциональна квадрату расстояния между ними.

\[F=G \frac{m_1 \cdot  m_2}{R_{12} ^2}\]

Таким образом, сила притяжения Земли к Солнцу $F_{\text{зс}}=G \frac{m_{\text{з}} \cdot  m_{\text{с}}}{R_{\text{зс}} ^2} ,$ а сила притяжения Марса к Солнцу $F_{\text{мс}}=G \frac{m_{\text{м}} \cdot  m_{\text{с}}}{R_{\text{мс}} ^2} .$ Чтобы ответить на вопрос задачи необходимо выразить отношение $\frac{R_{\text{мс}}}{R_{\text{зс}}} .$ Получаем $\frac{R_{\text{мс}}}{R_{\text{зс}}}=\sqrt{\frac{m_{\text{м}}}{m_{\text{з}}}\frac{F_{\text{зс}}}{F_{\text{мс}}}}=\sqrt{\frac1{10}\frac{22{,}5}{1}}=\sqrt{2{,}25}=1{,}5.$

Ответ: $1{,}5$.
}

% Source sdamgia-8933: https://phys-ege.sdamgia.ru/problem?id=8933
\BankTask{sdamgia-8933}{
Модуль ускорения свободного падения вблизи поверхности астероида равен $0{,}2\,\text{м/с}^{2}$. Чему будет равен модуль ускорения свободного падения вблизи поверхности другого астероида, объем которого в $8$ раз меньше? Оба астероида однородные, сферические и состоят из железа. (В метрах на секунду в квадрате.)
}{
\smallskip\noindent\textbf{Решение}\par
На тело вблизи астероида действует сила тяготения, которая связана с ускорением свободного падения:

\[F= \frac{GM_{\text{аст}1}m}{r_1 ^2} =mg_1.\]

Уменьшение объема астероида в $8$ раз приводит, во-первых, к уменьшению массы астероида в $8$ раз, а во-вторых, к уменьшению расстояния между центром астероида и телом вблизи его поверхности в $\sqrt[3]{8}=2$ раза. Таким образом:

\[g_2= \frac{GM_{\text{аст}2}}{r_2 ^2} = \frac{4}{8} \cdot  g_1=0{,}5g_1=0{,}1\,\text{м/с} ^2 .\]

Ответ: $0{,}1$.
}

% Source sdamgia-9494: https://phys-ege.sdamgia.ru/problem?id=9494
\BankTask{sdamgia-9494}{
Сила гравитационного взаимодействия небольших тел массами $m$ и $M$, находящихся на расстоянии $R_1=100\,\text{км}$ друг от друга, равна по модулю $F$. Сила гравитационного взаимодействия небольших тел массами $2m$ и $M$, находящихся на расстоянии $R_{2}$ друг от друга, равна по модулю $F/50$. На какую величину отличаются расстояния $R_{1}$ и $R_{2}$? (В километрах.)
}{
\smallskip\noindent\textbf{Решение}\par
Согласно закону всемирного тяготения сила притяжения между двумя материальными точками прямо пропорциональна произведению их масс и обратно пропорциональна квадрату расстояния между ними

\[F=G \frac{m_1 \cdot  m_2}{R ^2} .\]

Отношение сил гравитационного взаимодействия в первом и во втором случае $\frac{F_1}{F_2} =50= \frac{R_2 ^2}{2R_1 ^2} .$ 

Найдем отсюда $R_2$:

\[R_2=\sqrt{100R_1^2}=1000\,\text{км}.\]

Тогда разница $R_2 - R_1$ составляет $900\,\text{км}$.

Ответ: $900\,\text{км}$.
}

% Source sdamgia-19724: https://phys-ege.sdamgia.ru/problem?id=19724
\BankTask{sdamgia-19724}{
Однородная сферическая планета радиусом $R$ имеет плотность $\rho .$ Установите соответствие между физическими величинами и формулами, по которым эти величины можно рассчитать. 

\TaskMatchingColumns{Физическая величина}{А) ускорение свободного падения на поверхности планеты\par 
Б) первая космическая скорость для этой планеты}{Формула}{$1$) $2R\sqrt{\frac{\pi G\rho}{3}}$\par 
$2$) $\frac23R\sqrt{\pi G\rho}$\par 
$3$) $4 \pi  G\rho R$\par 
$4$) $\frac{4}{3} \pi  G\rho R$}

\begin{center}\small\renewcommand{\arraystretch}{1.25}\begin{tabular}{|p{0.450\linewidth}|p{0.450\linewidth}|}\hline
А & Б\\\hline
 & \\\hline
\end{tabular}\end{center}
}{
\smallskip\noindent\textbf{Решение}\par
А) Применим формулы ускорения свободного падения $g = \frac{GM}{R ^2} ,$ массы тела $M = \rho V,$ объема шара $V = \frac{4}{3} \pi  R ^3 .$ В результате получаем общую формулу:

\[g = \frac{4 \pi  G\rho R ^3}{3R ^2} = \frac{4 \pi  G\rho R}{3} .\]

Б) Применим формулу первой космической скорости $v=\sqrt{\frac{GM}{R}},$ и формулы массы и объема. Получим в результате: 

\[v=\sqrt{\frac{4\pi G\rho R^3}{3R}}=2R\sqrt{\frac{\pi G\rho}{3}}.\]

Ответ: $41$.
}

% Source sdamgia-3776: https://phys-ege.sdamgia.ru/problem?id=3776
\BankTask{sdamgia-3776}{
Вес тела на полюсе планеты, имеющей форму шара, на $16$ процентов превышает вес на экваторе. Чему равен период обращения планеты, если ее плотность $\rho = 0{,}7 \cdot  10 ^3\,\text{кг}/\text{м} ^3$? (В секундах, округлив до целого числа тысяч.)
}{
\smallskip\noindent\textbf{Решение}\par
\begin{center}\begin{tikzpicture}[book diagram,>=stealth]
\draw[thick](0,0)circle(1.6);\draw(-1.9,0)--(2.5,0);\draw(0,-1.9)--(0,3.2);
\draw(-.55,2.85)arc(190:530:.55 and .15);\node[right]at(.55,2.85){$\omega$};
\draw[fill=white](-.2,1.6)rectangle(.2,2.1);\draw[fill=white](1.6,-.2)rectangle(2.1,.2);
\draw[->,green!60!black,thick](0,1.6)--(0,2.6)node[right]{$N_1$};\draw[->,red,thick](0,1.6)--(0,.8)node[left]{$F$};
\draw[->,green!60!black,thick](1.6,0)--(2.5,0)node[above]{$N_2$};\draw[->,red,thick](1.6,0)--(.7,0)node[above]{$F$};
\draw[->,thick](0,0)--(-1.13,-1.13)node[midway,above left]{$R$};\node[left]at(-.2,2.45){$a=0$};
\draw[->,blue,thick](2.5,-1)--(1.65,-1)node[midway,below]{$\vec a$};
\end{tikzpicture}\end{center}Вес P тела — это сила, с которой тело давит на опору или растягивает подвес. Определим, чему равен вес некоторого тела на полюсе и на экваторе. Для определенности рассмотрим тело, лежащее на горизонтальной опоре (это предположение не ограничивает общность получаемых далее результатов). 

По третьему закону Ньютона сила действия равна силе противодействия, а потому, вес тела равен по величине силе реакции опоры $N$, действующей на тело. Следовательно, достаточно сравнить силы реакции опоры в двух случаях. 

Обозначим массу тела через $m$. Пусть радиус планеты через $R$, тогда масса планеты равна $M= \frac{4}{3} \pi  \rho R ^3 .$ Поскольку планета однородная и имеет форму сферы, согласно закону всемирного тяготения сила притяжения тела к планете в обеих точках одинакова и равна

\[F= \frac{GMm}{R ^2} = \frac{G ( \frac{4}{3} \pi  \rho R ^3 ) m}{R ^2} = \frac{4}{3} \pi  G\rho mR.\]

Рассмотрим случай, когда тело находится на полюсе. Второй закон Ньютона для него в проекции на радиальную ось приобретает вид:

\[\begin{gathered}N_1 - F=0.\end{gathered}\]

Когда же тело находится на экваторе, оно вращается вместе с планетой вокруг оси вращения, а значит, оно движется с центростремительным ускорением $a=\omega ^2 R,$ где $\omega$ — угловая скорость вращения планеты. Следовательно, для тела на экваторе второй закон Ньютона в проекции на радиальную ось приобретает вид: 

\[N_2 - F= - ma.\]

Согласно условию, $P_1=1{,}16P_2,$ а значит,

\[\begin{gathered}N_1=1{,}16N_2 \\\Longleftrightarrow  F=1{,}16 ( F - ma ) \Longleftrightarrow  1{,}16ma=0{,}16F \Longleftrightarrow \end{gathered}\]

\[1{,}16m\omega^2R=0{,}16\cdot\frac43\pi G\rho mR\Longleftrightarrow\omega=\sqrt{\frac{0{,}16}{1{,}16}\cdot\frac43\pi G\rho},\]

\[\omega=\sqrt{\frac{0{,}16}{1{,}16}\cdot\frac43\cdot3{,}14\cdot6{,}7\cdot10^{-11}\,\text{Н}\cdot\text{м}^2/\text{кг}^2\cdot0{,}7\cdot10^3\,\text{кг}/\text{м}^3}\approx1{,}65\cdot10^{-4}\,\text{с}^{-1}.\]

Отсюда для периода обращения планеты имеем:

\[T=\frac{2\pi}{\omega}=\frac{2\cdot3{,}14}{1{,}65\cdot10^{-4}\,\text{с}^{-1}}\approx3{,}8\cdot10^4\,\text{с}=38000\,\text{с}.\]

Ответ: $38000$.
}

% Source sdamgia-4439: https://phys-ege.sdamgia.ru/problem?id=4439
\BankTask{sdamgia-4439}{
Известно, что один оборот вокруг своей оси Венера совершает примерно за $243$ земных суток, а масса Венеры составляет $0{,}82$ от массы Земли. На орбиту какого радиуса надо вывести спутник Венеры, чтобы он все время «висел» над одной и той же точкой поверхности? Известно, что спутники Земли, «висящие» над одной и той же точкой поверхности, летают по орбите радиусом $R_{\text{з}}\approx42000\,\text{км}$.
}{
\smallskip\noindent\textbf{Решение}\par
При движении спутника по круговой орбите радиусом $R$ вокруг планеты центростремительное ускорение обеспечивается силой гравитационного притяжения спутника к планете, поэтому второй закон Ньютона записывается в виде $F_{\text{т}} = ma$ или 

\[m\omega ^2 R= \frac{GmM}{R ^2} ,\]

 где $m$ и $M$ — массы спутника и планеты, $G$ — гравитационная постоянная, $\omega = \frac{2 \pi }{T}$ — угловая скорость вращения спутника вокруг планеты, $T$ — период обращения спутника. Подставляя, получаем величину радиуса орбиты:

\[R=\sqrt[3]{\frac{GM}{\omega^2}}=\sqrt[3]{\frac{GMT^2}{4\pi^2}}.\]

Из полученной формулы найдем отношение радиусов орбит спутников Венеры и Земли: 

\[\frac{R_{\text{В}}}{R_{\text{З}}}=\sqrt[3]{\frac{GM_{\text{В}}T_{\text{В}}^2}{4\pi^2}}:\sqrt[3]{\frac{GM_{\text{З}}T_{\text{З}}^2}{4\pi^2}}=\sqrt[3]{\frac{M_{\text{В}}}{M_{\text{З}}}\cdot\frac{T_{\text{В}}^2}{T_{\text{З}}^2}}.\]

Период обращения спутника вокруг Земли $T_{\text{З}} = 1$ сутки, поэтому, подставляя данные из условия, находим:

\[R_{\text{В}}=\sqrt[3]{0{,}82\cdot243^2}R_{\text{З}}\approx1531000\,\text{км}.\]

Ответ: $R_{\text{В}}=R_{\text{З}}(0{,}82\cdot243^2)^{1/3}\approx1531000\,\text{км}.$

Примечание Д. Д. Гущина.

Внимательный читатель заметит, что найденный радиус орбиты превышает радиус сферы Хилла, равный для Венеры одному миллиону километров. Иными словами, притяжения Венеры не хватает, чтобы удерживать настолько удаленный от нее спутник. Поэтому со временем он покинет орбиту Венеры и станет независимо от Венеры вращаться вокруг Солнца. Чтобы этого не произошло, необходимо (пока позволяют запасы топлива) корректировать орбиту спутника двигателями.

Пытливый читатель может заинтересоваться вопросом о том, каков радиус сферы Хилла для Земли. Он равен $1{,}47\,\text{млн\,км}$. С удовольствием отметим, что Луна удалена от Земли на $0{,}384\,\text{млн\,км}$, поэтому Луна не покинет орбиту Земли. Подробности об этом можно прочитать в англоязычной Википедии.
}

% Source sdamgia-31899: https://phys-ege.sdamgia.ru/problem?id=31899
\BankTask{sdamgia-31899}{
$22$ октября 2022 года с космодрома «Восточный» ракетой из семейства «Союз» были запущены три спутника связи «Гонец-М» и первый спутник «Скиф-Д» из проекта «Сфера» по обеспечению доступа к широкополосному интернету на всей территории России. Этот спутник вышел на круговую орбиту высотой $h=8170\,\text{км}$ над поверхностью Земли с полярным наклоном $90^\circ$ (плоскость такой орбиты проходит через оба географических полюса Земли), и через некоторое время пролетел с юга на север над центром Москвы. Через сколько оборотов вокруг Земли он снова окажется над Москвой, двигаясь в том же направлении с юга на север? Радиус Земли $R_{\text{з}}=6370\,\text{км}$.
}{
\smallskip\noindent\textbf{Решение}\par
$1$. Плоскость орбиты спутника сохраняет свое положение в пространстве, а Земля вертится вокруг оси, и точки на ее поверхности пересекают эту плоскость по два раза за период обращения спутника, причем нас интересует только одно пересечение, когда спутник пролетает над Москвой в направлении с юга на север.

$2$. Чтобы спутник снова пролетел над Москвой в том же направлении, очевидно, необходимо выполнение соотношения $mT_{\text{з}} = nT_{\text{сп}},$ где $m$ и $n$ — целые числа, а $T_{\text{з}}=24\,\text{часа}=1440\,\text{мин}$ и $T_{\text{сп}}$ — периоды обращения Земли вокруг своей оси и спутника по круговой полярной орбите вокруг Земли.

$3$. Для нахождения $T_{\text{сп}}= \frac{2 \pi  r}{v _1} ,$ где $r=R_{\text{з}}+h$, а $v _1$— скорость спутника массой $m_{\text{с}}$, запишем для него второй закон Ньютона с учетом выражений для силы тяжести Земли $F= \frac{GMm_{\text{с}}}{r ^2}$ и центростремительного ускорения $a= \frac{v _1 ^2}{r}$:

\[\frac{GMm_{\text{с}}}{r ^2} = \frac{m_c v _1 ^2}{r} ,\]

 откуда $v_1=\sqrt{\frac{GM}{r}},$ и поскольку $g_0= \frac{GM}{R ^2 _{\text{з}}} ,$ то:

\[T_{\text{сп}}=2\pi r\sqrt{\frac{r}{GM}}=\frac{2\pi}{R_{\text{з}}}\sqrt{\frac{(R_{\text{з}}+h)^3}{g_0}}.\]

 Здесь $g_{0}$ и $g_1$ — ускорения свободного падения на поверхности Земли и на высоте $h$ над ней.

$4$. Подставляя численные данные из условия и таблиц, получаем:

\[T_{\text{сп}}=\frac{2\cdot3{,}14}{6370000}\sqrt{\frac{(6370000+8170000)^3}{10}}\approx17285\,\text{с}\approx288{,}08\,\text{мин}.\]

$5$. Таким образом:

\[\frac{n}{m} = \frac{T_{\text{з}}}{T_{\text{сп}}} \approx \frac{1440}{288{,}08} \approx 4{,}999\approx 5,\]

 то есть за одни сутки спутник совершит $5$ оборотов вокруг Земли и снова окажется над Москвой.

Ответ: $5$ оборотов.
}


% Source AC8852: https://ege.fipi.ru/bank/questions.php?search=1&proj=BA1F39653304A5B041B656915DC36B38&qid=AC8852
\BankTask{AC8852}{
Расстояние от искусственного спутника до поверхности Земли равно двум радиусам Земли. Во сколько раз увеличится сила притяжения спутника к Земле, если расстояние от него до поверхности Земли станет равным одному радиусу Земли?
}{
\smallskip\noindent\textbf{Решение от GPT}\par
Первоначально расстояние до центра Земли $r_1=3R$, затем $r_2=2R$. По закону всемирного тяготения $F=GMm/r^2$, поэтому\[\frac{F_2}{F_1}=\left(\frac{r_1}{r_2}\right)^2=\left(\frac{3}{2}\right)^2=2{,}25.\] Ответ: в $2{,}25$ раза.
}

% Source 9F5549: https://ege.fipi.ru/bank/questions.php?search=1&proj=BA1F39653304A5B041B656915DC36B38&qid=9F5549
\BankTask{9F5549}{
Два маленьких шарика массой $m$ каждый находятся на расстоянии $40\,\text{см}$ друг от друга. Каково расстояние между шариками вдвое большей массы, если модуль сил гравитационного взаимодействия между ними такой же, как и между первыми двумя шариками? Результат выразите в $\text{см}$.
}{
\smallskip\noindent\textbf{Решение от GPT}\par
$Gm^2/r_1^2=G(2m)^2/r_2^2$, откуда $r_2=2r_1=80\,\text{см}$.
}

% Source 3FB779: https://ege.fipi.ru/bank/questions.php?search=1&proj=BA1F39653304A5B041B656915DC36B38&qid=3FB779
\BankTask{3FB779}{
Две планеты обращаются по круговым орбитам вокруг звезды. Модули сил притяжения планет к звезде одинаковы. У первой из планет радиус орбиты
в шесть раз больше, чем у второй. Каково отношение $\frac{m_1}{m_2}$ массы первой планеты к массе второй планеты?
}{
\smallskip\noindent\textbf{Решение от GPT}\par
При одинаковых силах $m_1/r_1^2=m_2/r_2^2$. Следовательно, $m_1/m_2=(r_1/r_2)^2=6^2=36$.
}

% Source 41BBCF: https://ege.fipi.ru/bank/questions.php?search=1&proj=BA1F39653304A5B041B656915DC36B38&qid=41BBCF
\BankTask{41BBCF}{
Два маленьких шарика массой $m$ каждый находятся на расстоянии $r$ друг от друга и притягиваются друг к другу с силой $8\,\text{нН}$. Какова сила гравитационного притяжения двух других шариков, если масса одного $2m$, масса другого $\frac{m}{2}$, а расстояние между их центрами $2r$? Результат выразите в $\text{нН}$.
}{
\smallskip\noindent\textbf{Решение от GPT}\par
$F_2/F_1=(2m\cdot m/2)/m^2\cdot(r/(2r))^2=1/4$. Поэтому $F_2=8/4=2\,\text{нН}$.
}

% Source BA66EA: https://ege.fipi.ru/bank/questions.php?search=1&proj=BA1F39653304A5B041B656915DC36B38&qid=BA66EA
\BankTask{BA66EA}{
У поверхности Луны на космонавта действует сила тяготения $120\,\text{Н}$. Какая сила тяготения действует со стороны Луны на того же космонавта
в космическом корабле, движущемся по круговой орбите вокруг Луны
на расстоянии двух лунных радиусов от её центра? Результат выразите в $\text{Н}$.
}{
\smallskip\noindent\textbf{Решение от GPT}\par
$F\sim r^{-2}$, поэтому на расстоянии $2R$ от центра Луны $F_2=F_1/4=120/4=30\,\text{Н}$.
}

% Source A129A4: https://ege.fipi.ru/bank/questions.php?search=1&proj=BA1F39653304A5B041B656915DC36B38&qid=A129A4
\BankTask{A129A4}{
Два одинаковых маленьких шарика массой $m$ каждый, расстояние между центрами которых равно $r$, притягиваются друг к другу с силами, равными по модулю $0{,}2\,\text{пН}$. Каков модуль сил гравитационного притяжения двух других шариков, если масса каждого из них $2m$, а расстояние между их центрами $2r$? Результат выразите в $\text{пН}$.
}{
\smallskip\noindent\textbf{Решение от GPT}\par
$F_2/F_1=(2m)^2/m^2\cdot(r/(2r))^2=1$. Сила не изменится: $F_2=0{,}2\,\text{пН}$.
}

\end{document}
